\appendix


\chapter{Código: Curvas de Desempenho}
\label{anexo:codigo_pythonb}

O código a seguir apresenta a implementação do script em Python, utilizado para o cálculo da variação dos parâmetros operacionais (altura manométrica e potência hidráulica) com base nas Leis de Afinidade para turbomáquinas, bem como para a geração gráfica das curvas de desempenho em função da rotação do eixo.
% Configuração para aceitar acentos em português dentro do código Python
\lstset{literate=
    {á}{{\'a}}1 {é}{{\'e}}1 {í}{{\'i}}1 {ó}{{\'o}}1 {ú}{{\'u}}1
    {Á}{{\'A}}1 {É}{{\'E}}1 {Í}{{\'I}}1 {Ó}{{\'O}}1 {Ú}{{\'U}}1
    {à}{{\`a}}1 {è}{{\`e}}1 {ì}{{\`i}}1 {ò}{{\`o}}1 {ù}{{\`u}}1
    {À}{{\`A}}1 {È}{{\'E}}1 {Ì}{{\`I}}1 {Ò}{{\`O}}1 {Ù}{{\`U}}1
    {ä}{{\"a}}1 {ë}{{\"e}}1 {ï}{{\"i}}1 {ö}{{\"o}}1 {ü}{{\"u}}1
    {Ä}{{\"A}}1 {Ë}{{\"E}}1 {Ï}{{\"I}}1 {Ö}{{\"O}}1 {Ü}{{\"U}}1
    {â}{{\^a}}1 {ê}{{\^e}}1 {î}{{\^i}}1 {ô}{{\^o}}1 {û}{{\^u}}1
    {Â}{{\^A}}1 {Ê}{{\^E}}1 {Î}{{\^I}}1 {Ô}{{\^O}}1 {Û}{{\^U}}1
    {ã}{{\~a}}1 {ẽ}{{\~e}}1 {ĩ}{{\~i}}1 {õ}{{\~o}}1 {ũ}{{\~u}}1
    {Ã}{{\~A}}1 {Ẽ}{{\~E}}1 {Ĩ}{{\~I}}1 {Õ}{{\~O}}1 {Ũ}{{\~U}}1
    {œ}{{\oe}}1 {Œ}{{\OE}}1 {æ}{{\ae}}1 {Æ}{{\AE}}1 {ß}{{\ss}}1
    {ű}{{\H{u}}}1 {Ű}{{\H{U}}}1 {ő}{{\H{o}}}1 {Ő}{{\H{O}}}1
    {ç}{{\c c}}1 {Ç}{{\c C}}1 {ø}{{\o}}1 {å}{{\r a}}1 {Å}{{\r A}}1
    {€}{{\euro}}1 {£}{{\pounds}}1 {«}{{\guillemotleft}}1
    {»}{{\guillemotright}}1 {ñ}{{\~n}}1 {Ñ}{{\~N}}1 {¿}{{?`}}1 {¡}{{!`}}1
}

\begin{lstlisting}[language=Python, caption={Script de simulação cinemática da Lewis VL750}]
# %%

if __name__ == "__main__":

    # Premissas do Projeto (Carga de ~10,11 m)
    D = 0.93
    d = 0.60
    Q_base = 2.2222
    alpha_4 = 90.0
    beta_5 = 8.87  # Ajustado para bater exatamente com a altura manométrica de 10.11 m

    # Criando o analisador com os dados do problema
    analyzer = TurbomachineAnalyzer(
        D=D,
        d=d,
        N=975,
        Q=Q_base,
        alpha_4=alpha_4,
        beta_5=beta_5
    )

    # Pegando resultados numéricos
    print("GEOMETRIA:")
    for k, v in analyzer.get_geometry().items():
        print(f"  {k}: {v:.4f}" if isinstance(v, float) else f"  {k}: {v}")

    print("\nSEÇÃO 4 (Entrada):")
    for k, v in analyzer.get_section4().items():
        print(f"  {k}: {v:.4f}" if isinstance(v, float) else f"  {k}: {v}")

    print("\nSEÇÃO 5 (Saída):")
    for k, v in analyzer.get_section5().items():
        print(f"  {k}: {v:.4f}" if isinstance(v, float) else f"  {k}: {v}")

    print("\nTRANSFERÊNCIA DE ENERGIA (Euler):")
    for k, v in analyzer.get_energy_transfer().items():
        print(f"  {k}: {v:.4f}")

    # Plotar triângulos
    analyzer.plot_velocity_triangles(show=True)
\end{lstlisting}